\newpage
\section{Lookup soundness calculation}

We consider logUP, the lookup argument introduced in~\cite{habock2022logup}. 

We have a table $\mathbf{T}\in\mathbb{F}^{T\times S}$ and tables $\mathbf{L}^1, \ldots, \mathbf{L}^K$ in $\mathbb{F}^{L\times S}$ that are lookups of it, that is, $\forall \ i\in[L] \ \exists \ j\in[T] \ s.t. \ \mathbf L_i=\mathbf T_j$.

For the sake of the integration and in order to be consistent with the implementation parameters, we consider separately the cases when the lookup argument is proven using a univariate or multivariate IOP. We also consider a different case when $S=1$.

In the code, the parameter controlling the aggregation is \texttt{multilinear\_fingerprint}. If it is not specified, then we assume it follows the lookup type:
\begin{itemize}
	\item univariate lookup $\Rightarrow$ univariate aggregation,
	\item multivariate lookup $\Rightarrow$ multivariate aggregation.
\end{itemize}
Only when \texttt{multilinear\_fingerprint} is explicitly specified do we override this default.

The logUP argument proves lookup relations by proving that given a vector $\vec m\in\mathbb{F}^{|\mathbf T|}$ such that $m_j\in\mathbb{Z}$, the lookup relation holds if and only if the following equation is satisfied.

$$\sum_{\vec l\in\mathbf{L}} \frac{1}{X-  \vec l}=\sum_{\vec t \in \mathbf{T}} \frac{m_{\vec t}}{X- \vec t}$$

The expression above is equivalent to considering only one table $\mathbf L$ and having each element taken a positive or negative value depending on whether it is sent from the current trace or received from some previous computation. In that case, we just set $\mathbf{L}=\mathbf 0$ and have
$$\sum_{\vec t \in \mathbf{T}} \frac{m_{\vec t}}{X- \vec t}=0$$

We split the logUP protocol in two steps:

\begin{enumerate}
	\item Compute some polynomial expression of the equation above
	\item Prove consistency with $\mathbf T$ and $\mathbf L$	
\end{enumerate}



\subsection{Univariate case}

In this case, we always assume that proving the well formation of the sums is deferred to AIR constraints, and that the additional constraints this causes are already taken into account in the size of the AIR trace for the soundness of the proof of correct execution of it.

If \texttt{multilinear\_fingerprint} is not specified, then in this section we assume it is \texttt{false}, i.e. the aggregation is univariate.

We consider the following parameters:
\begin{itemize}
	\item $F$ = size of the extension field $\mathbb{F}$
	\item $H_L$ = size of the set used to interpolate $\mathbf{L}$
	\item $H_T$ = size of the set used to interpolate $\mathbf{T}$
	\item $H = H_L + H_T$
	\item $S$ = number of columns of $\mathbf{T}$ and $\mathbf{L}$
	\item $K$ = number of lookups performed on $\mathbf{T}$
\end{itemize}


The expression 

$$\sum_{i=1}^L \frac{1}{X- \sum_{s=1}^S \mathbf L_{is}Y^{s-1}}=\sum_{j=1}^T \frac{m_j}{X-\sum_{s=1}^S \mathbf T_{js}Y^{s-1}}$$

is replaced, upon receiving $\alpha, \beta$ uniformly sampled from $\mathbb{F}$, by the polynomial equality  

$$\sum_{i=1}^L \frac{1}{\alpha- \sum_{s=1}^S \mathbf L_{is}\beta^{s-1}}\lambda_i(X)=\sum_{j=1}^T \frac{m_j}{\alpha-\sum_{s=1}^S \mathbf T_{js}\beta^{s-1}}\mu_j(X)$$

$\lambda_i(X), \mu_j(X)$ being the Lagrange interpolation polynomials of some domains of size $L$ and $T$, respectively.


The probability that a prover gets a randomness that verifies without $\mathbf L$ being a lookup on $\mathbf{T}$ is bounded by $\epsilon_{\mathsf{sum}}$ as follows:


\subsubsection{Multi-column and aggregation:}

Let $R$ be the contribution of the column aggregation. In the default univariate setting of this section, we take $R=S$. If the implementation explicitly overrides the aggregation to be multivariate, then we instead take $R=\log(S)$. Thus:
\[
R =
\begin{cases}
S & \text{if the aggregation is univariate,}\\
\log(S) & \text{if the aggregation is multivariate.}
\end{cases}
\]

Then the lookup soundness error is bounded by
\[
\epsilon_{\mathsf{sum}}\leq \frac{K H R}{F}.
\]

For the single-column case, we take $R=1$.

\subsubsection{Dummy Variables}

If $\mathbf T, \mathbf L$ are smaller than the interpolation set and the tables are filled with dummy variables, those variables should be taken into account in the sizes of $H_L$ and $H_T$. 

\subsection{Multivariate}

If \texttt{multilinear\_fingerprint} is not specified, then in this section we assume it is \texttt{true}, i.e. the aggregation is multivariate.

We consider the following parameters:
\begin{itemize}
	\item $F$ = size of the extension field $\mathbb{F}$
	\item $H_L$ = size of the alphabet used to interpolate $\mathbf{L}$
	\item $H_T$ = size of the alphabet used to interpolate $\mathbf{T}$
	\item $H = H_L + H_T$
	\item $S$ = number of columns of $\mathbf{T}$ and $\mathbf{L}$
	\item $K$ = number of lookups performed on $\mathbf{T}$
\end{itemize}

The expression 

$$\sum_{i=1}^L \frac{1}{X- \sum_{s=1}^S \mathbf L_{is}Y^{s-1}}=\sum_{j=1}^T \frac{m_j}{X-\sum_{s=1}^S \mathbf T_{js}Y^{s-1}}$$

is replaced, upon receiving $\alpha$ uniformly sampled from $\mathbb{F}$ and $\vec \beta$ uniformly sampled from $\mathbb{F}^{\log S}$, by the polynomial equation:  

$$\sum_{\vec x\in H} \frac{1}{\alpha- L(\vec x, \vec \beta)}=\sum_{\vec x\in H} \frac{m(\vec x)}{\alpha-T(\vec x, \vec \beta)}$$


To convert the equation above into a polynomial one, both sides are multiplied by 

$$\prod_{\vec x\in H}(\alpha - L(\vec x, \vec \beta))\prod_{\vec x\in H}(\alpha - T(\vec x, \vec \beta))$$

\subsubsection{Multi-column and aggregation:}

Let $R$ be the contribution of the column aggregation. In the default multivariate setting of this section, we take $R=\log(S)$. If the implementation explicitly overrides the aggregation to be univariate, then we instead take $R=S$. Thus:
\[
R =
\begin{cases}
S & \text{if the aggregation is univariate,}\\
\log(S) & \text{if the aggregation is multivariate.}
\end{cases}
\]

Then the lookup soundness error is bounded by
\[
\epsilon_{\mathsf{sum}}\leq \frac{K H R}{F}.
\]

For the single-column case, we take $R=1$.

To prove consistency with $\mathbf L$ and $\mathbf T$, we also account for the
GKR soundness term. Writing $\epsilon_{\mathsf{gkr}}$ for that contribution, we
use
\[
\epsilon_{\mathsf{gkr}}
\leq \frac{1}{2}\frac{(\log(H)+\log(K))(3(\log(H)+\log(K))+1)}{F},
\]
which matches the implementation with $H = H_L + H_T$.

\subsection{Grinding}

For both cases we capture the option to decrease the soundness error of the lookup round through grinding. We consider a parameter $\mathsf{grinding\_bits\_lookup}$ denoting the rounds of proof of work, and compute the soundness error as
\[
\epsilon=
\begin{cases}
\epsilon_{\mathsf{sum}}2^{-\mathsf{grinding\_bits\_lookup}} & \text{in the univariate case,}\\
(\epsilon_{\mathsf{sum}}+\epsilon_{\mathsf{gkr}}+\epsilon_{\mathsf{reduction}})2^{-\mathsf{grinding\_bits\_lookup}} & \text{in the multivariate case.}
\end{cases}
\]
